\subsection{SATURATED SUBSETS}

Let \(\mathbb{R}\{x, y\}\) be an equivalence relation on a set \(\mathbf{E}\) , and let \(\mathbf{A}\) be a subset of \(\mathbf{E}\) . A is said to be saturated with respect to \(\mathbb{R}\) if the relation \(x \in \mathbf{A}\) is compatible (with respect to \(x\) ) with \(\mathbb{R}\{x, y\}\) ; or, equivalently, if for each \(x \in \mathbf{A}\) the equivalence class of \(x\) is contained in \(\mathbf{A}\) . In other words, a set is saturated with respect to \(\mathbb{R}\) if and only if it is the union of a set of equivalence classes with respect to \(\mathbb{R}\) .

Let \(f\) be the canonical mapping of \(\mathbf{E}\) onto \(\mathbf{E} / \mathbf{R}\) . If \(\mathbf{A}\) is saturated with respect to \(\mathbf{R}\) , then the equivalence class of each \(x \in \mathbf{A}\) , which is equal to \(\overline{f}^1 \langle \{f(x)\} \rangle\) , is contained in \(\mathbf{A}\) , and hence \(\overline{f}^1 \langle f\langle \mathbf{A}\rangle \rangle \subset \mathbf{A}\) ; since in any case we have \(\mathbf{A} \subset \overline{f}^1 \langle f\langle \mathbf{A}\rangle \rangle\) , it follows that \(\overline{f}^1 \langle f\langle \mathbf{A}\rangle \rangle = \mathbf{A}\) . Conversely, if \(\mathbf{A} = \overline{f}^1 \langle f\langle \mathbf{A}\rangle \rangle\) , then for every \(x \in \mathbf{A}\) the equivalence class \(\mathbf{K} = f(x)\) of \(x\) with respect to \(\mathbf{R}\) is an element of \(f\langle \mathbf{A}\rangle\) ; and since \(\mathbf{K} = \overline{f}^1 \langle \{\mathbf{K}\}\rangle\) , we have \(\mathbf{K} \subset \overline{f}^1 \langle f\langle \mathbf{A}\rangle \rangle = \mathbf{A}\) . Thus the subsets of \(\mathbf{E}\) which are saturated with respect to \(\mathbf{R}\) are precisely the subsets \(\mathbf{A}\) of \(\mathbf{E}\) such that \(\mathbf{A} = \overline{f}^1 \langle f\langle \mathbf{A}\rangle\rangle\) . Equivalently, they are the subsets of \(\mathbf{E}\) of the form \(\overline{f}^1\langle\mathbf{B}\rangle\) , where \(\mathbf{B} \subset E / R\) ; for the relation \(A = \overline{f}^1\langle B\rangle\) implies \(B = f\langle A\rangle\) , hence \(A = \overline{f}^1\langle f\langle A\rangle\rangle\) .

If \((\mathbf{X}_t)_{t\in I}\) is a family of saturated subsets of \(E\) , then the sets \(\bigcup_{t\in I} X_t\) and \(\bigcap_{t\in I} X_t\) are saturated (§4, Propositions 3 and 4). If \(A\) is a saturated subset of \(E\) , then so is \(\complement_E A\) (§4, Proposition 6).

Let A now be an arbitrary subset of E. Then the set \(\overline{f}^1\langle f\langle \mathrm{A}\rangle \rangle\) contains A and is saturated. Conversely, if \(\mathbf{A}'\) is a saturated subset of E which contains A, we have \(f\langle \mathrm{A}'\rangle \supset f\langle \mathrm{A}\rangle\) , so that

\[
\mathrm{A} ^ {\prime} = \overline {{f}} ^ {1} \langle f \langle \mathrm{A} ^ {\prime} \rangle \rangle \supset \overline {{f}} ^ {1} \langle f \langle \mathrm{A} \rangle \rangle .
\]

Hence we may say that \(\overline{f}^1\langle f\langle A\rangle \rangle\) is the ``smallest'' saturated subset of E which contains A (cf.~Chapter III); this set is called the saturation of A with respect to the relation R. It is immediately seen that the saturation of A is the union of the equivalence classes of the elements of A. If \((\mathbf{X}_{i})_{i\in I}\) is a family of subsets of E and if \(A_{i}\) is the saturation of \(X_{i}\) with respect to R, then the saturation of \(\bigcup_{i\in I}X_{i}\) is \(\bigcup_{i\in I}A_{i}\) (§4, Proposition 3).

